% Figure: a single rigid body (a bent bracket) with its complete planar
% free-body diagram: applied load, pin reaction (two components), roller
% reaction (one component normal to the surface), and the moment arm used
% in the moment sum about the pin.
\begin{figure}[htbp]
\centering
\begin{tikzpicture}[>=latex, scale=1.0]
  % the body: an L-bracket
  \fill[engblue!12, draw=engblue, thick]
    (0,0) -- (4,0) -- (4,0.5) -- (0.5,0.5) -- (0.5,2.5) -- (0,2.5) -- cycle;
  % pin support at origin (triangle + hatch)
  \fill[enggray] (0,0) circle (2pt);
  \draw[thick] (-0.35,-0.5) -- (0,0) -- (0.35,-0.5) -- cycle;
  \foreach \x in {-0.35,-0.2,...,0.4} {
    \draw[enggray, thin] (\x,-0.5) -- (\x-0.12,-0.7);
  }
  \node[anchor=north east] at (-0.1,-0.05) {$A$};
  % roller support under the right end
  \draw[thick] (4,-0.05) circle (0.12);
  \draw[thick] (3.7,-0.17) -- (4.3,-0.17);
  \foreach \x in {3.7,3.85,...,4.3} {
    \draw[enggray, thin] (\x,-0.17) -- (\x-0.12,-0.37);
  }
  \node[anchor=north] at (4,-0.42) {$B$};
  % applied load at the top of the vertical arm
  \draw[->, very thick, engred] (0.25,2.5) -- (0.25,1.4);
  \node[engred, anchor=west] at (0.3,2.0) {$P$};
  % pin reaction components at A
  \draw[->, thick, englime!70!black] (0,0) -- (1.1,0);
  \node[anchor=south] at (0.9,0.05) {$A_x$};
  \draw[->, thick, englime!70!black] (0,0) -- (0,1.1);
  \node[anchor=east] at (-0.05,0.9) {$A_y$};
  % roller reaction at B (vertical)
  \draw[->, thick, englime!70!black] (4,0) -- (4,1.1);
  \node[anchor=west] at (4.05,0.9) {$B_y$};
  % moment-arm annotation
  \draw[<->, enggray, thin] (0,-0.95) -- (4,-0.95);
  \node[enggray, anchor=north] at (2,-0.98) {$L$};
\end{tikzpicture}
\caption[Free-body diagram of a planar rigid body]{The free-body
diagram of a planar rigid body resolves every contact into the
reactions the support can transmit. A pin (left, support $A$) transmits
two force components $A_x$ and $A_y$ and no moment, so the body may
still rotate about the pin until a second support restrains it. A roller
(right, support $B$) transmits one force $B_y$ normal to its running
surface and nothing along it. The applied load $P$ and three reaction
components are four unknowns; the body in two dimensions supplies three
equations, so this configuration is statically determinate only because
the roller removes one component a second pin would have added. Taking
moments about $A$ eliminates $A_x$ and $A_y$ at a stroke and solves for
$B_y$ directly, which is the tactical reason to sum moments about a pin.}
\label{fig:vol03ch02:fbd-conventions}
\end{figure}
